\originalpreface
理论力学用坐标、变分和相空间描述运动. 本册面向已学牛顿力学、微积分、向量与矩阵运算的读者, 将8.223的分析力学基础和8.09的主体内容按学科重组. 正文从约束与拉格朗日方程出发, 推导对称守恒、中心力、小振动、刚体、哈密顿方程、正则变换、Hamilton--Jacobi方法和作用角变量, 再讨论摄动、连续介质和非线性动力学的入门内容. 新增解释、推导与例题属于中文整理者的补充, 不是逐页翻译的官方讲稿.

主源是 Iain W. Stewart 的 MIT 8.09 Classical Mechanics III, Fall 2014. 实际合本有199页, 版权页标2016, 编辑为 Francisco Machado、Thanawuth Thanathibodee、Prashanth S. Venkataram. 基础来源是 Matthew Evans 的8.223 Classical Mechanics II, January IAP 2017; 该课讲稿由听课学生打字整理, 学生个人署名未在课程页逐项列出. 原文存在公式错序、符号约定和个别论断不严密处, 中文稿以可核查的推导为准, 具体修正随题说明.

本册收录8.223四份作业的40道原题及中文解答, 包括原题的选做部分; 对有官方解答的题目标为“官方解答中文改写”, 另行补足或修正之处明确说明. 8.09作业8--10的十六题、作业1--7逐题核查后选择的二十题与投石机分析项目给出AI独立解答或数值复现, 不冒充MIT答案. 无阻尼受迫摆的一项周期窗口要求仍未完成, 具体子问状态随题说明. 项目的实物制作、演示视频和开放优化没有唯一纸面标准答案, 本册提供模型及验收方法, 不伪造实验成果. 8.09七份混合作业中明确商业教材来源题仍排除, 原PDF不打包; 选择的课程题以实际题面和有限权利核查为依据, 不将其视作一切潜在来源的绝对许可. 8.09考试确实列入课历, 但官方页面没有可取得的期中、期末试卷或解答, 本册不补造这些题面.

MIT OCW原件及本册相关改编按CC BY-NC-SA 4.0提供, 原资源另有权利声明时以其声明为准. 本册为独立中文学习整理, 不代表MIT官方出版、中文译本或认可. 未下载或打包Goldstein、Landau--Lifshitz等商业教材; 教材名仅用于交代原课阅读要求. 原件与逐文件URL、标题、作者、学期、校验值、页数及例外另存于项目来源清单.

\par\noindent{\bfseries 课程资料与课历}\par
8.09要求先修8.223且成绩至少C, 推荐8.03振动与波. 每周两次讲课, 每次1.5小时, 另一次1小时习题课. 评分为作业35\%, 两小时期中25\%, 三小时期末40\%. 官方大纲依次列分析力学与约束、刚体与转动、小振动、正则变换与Hamilton--Jacobi、摄动、流体、混沌. 公开课历如下, 交作业时间均为指定讲次前一天.
\begin{longtable}{p{.13\textwidth}p{.47\textwidth}p{.31\textwidth}}
\toprule 讲次 & 内容 & 考核节点\\\midrule
1--4 & 分析力学复习 & 作业1: 第4讲前\\
5--7 & 刚体动力学 & 作业2: 第6讲前\\
8--9 & 振动与振荡 & 作业3: 第8讲前\\
9--13 & 正则变换、HJ与作用角 & 作业4: 第10讲前; 作业5: 第13讲前\\
13--14 & 摄动论 & 作业6: 第14讲前\\
15--16 & 流体及期中复习 & 两小时期中\\
17--20 & 流体续讲 & 作业7: 第20讲前\\
21--26 & 混沌 & 作业8、9、10: 第21、23、26讲前\\
期末 & 综合考核 & 三小时期末\\\bottomrule
\end{longtable}
8.223先修8.01与18.02, 每周两次1小时讲课和一次1小时习题课. 大纲安排为第1--4讲拉格朗日力学, 5--6讲守恒量, 7--10讲轨道与散射, 11--12讲振动, 13--14讲特殊势, 15--19讲哈密顿力学. 评分为四份作业70\%, 项目30\%; 项目页面明确本课无考试, 独立分析为主, 指定实验或数值部分另有合作或使用计算机的例外.

8.223没有单独公开逐日课历, 本册保留上述讲次安排, 不推造日历日期. 已取得第1--6、8、10--18讲中的16份实际PDF, 第7、9、19讲未取得. 目录说第12讲未提供, 但独立资源有10页可读PDF: 网页标题为Forced Oscillations, 正文标题为Damped Oscillators. 本册按实际正文归入阻尼振动, 并保留这一来源差异. 项目页面标题标2016, 所在课程标2017 IAP, 也不将两者擅自改成一致.
